% GENERATED FILE. Edit canonical lessons, then run npm run generate:books.
% canonical-source-hash: fnv1a64:89c57276662dc4d0
% canonical-lessons: ML-C209-kariveppila, ML-C209-katuku, ML-C209-masala, ML-C209-cemmin, ML-C209-sandesam

\chapter{Curry leaves, Mustard seeds, Spice mix, Prawn, Message}
\label{ch:ml-curry-leaves-mustard-seeds-spice-mix-prawn-message}
\hlchaptermodality{209}

\begin{chapteropening}
\textbf{By the end of this chapter:} \emph{I can say curry leaves, mustard seeds, spice mix, a prawn and a message.}

Complete the last lesson of chapter 209: I can say curry leaves, mustard seeds, spice mix, a prawn and a message.
\end{chapteropening}

\section[\texorpdfstring{\ml{കറിവേപ്പില} (kaṟivēppila)}{കറിവേപ്പില (kaṟivēppila)}]{\ml{കറിവേപ്പില} (kaṟivēppila) — curry leaves}

\label{lesson:ML-C209-kariveppila}



\begin{quote}
Before the new one: say the Malayalam for butter, then the Malayalam for honey.
\end{quote}

\subsection*{You'll want to know: \ml{കറിവേപ്പില}}
\textbf{\ml{കറിവേപ്പില}} — \emph{kaṟivēppila} — \textquotedblleft{}curry leaves\textquotedblright{}.

\begin{grammarlens}[title={What you've built}]
One more word for this chapter.
\end{grammarlens}

\subsection*{Your turn}
\begin{itemize}
\raggedright
  \item \emph{Say it:} \emph{kaṟivēppila}
  \item \emph{Say it:} \emph{kaṟivēppila}, once more
  \item \emph{Say it:} say \emph{kaṟivēppila}
  \item \emph{Recall:} say the Malayalam for butter, then the Malayalam for honey, then say \emph{kaṟivēppila} again
\end{itemize}

\begin{tcolorbox}[breakable,colback=teal!4,colframe=teal!35!black,title={Before you move on}]
What does \ml{കറിവേപ്പില} mean? (\textquotedblleft{}Curry leaves\textquotedblright{}.) Say it once more.
\end{tcolorbox}
\section[\texorpdfstring{\ml{കടുക്} (kaṭukŭ)}{കടുക് (kaṭukŭ)}]{\ml{കടുക്} (kaṭukŭ) — mustard seeds}

\label{lesson:ML-C209-katuku}



\begin{quote}
Before the new one: say the Malayalam for honey, then the Malayalam for curry leaves.
\end{quote}

\subsection*{You'll want to know: \ml{കടുക്}}
\textbf{\ml{കടുക്}} — \emph{kaṭukŭ} — \textquotedblleft{}mustard seeds\textquotedblright{}.

\begin{grammarlens}[title={What you've built}]
One more word for this chapter.
\end{grammarlens}

\subsection*{Your turn}
\begin{itemize}
\raggedright
  \item \emph{Say it:} \emph{kaṭukŭ}
  \item \emph{Say it:} \emph{kaṭukŭ}, once more
  \item \emph{Say it:} say \emph{kaṭukŭ}
  \item \emph{Recall:} say the Malayalam for honey, then the Malayalam for curry leaves, then say \emph{kaṭukŭ} again
\end{itemize}

\begin{tcolorbox}[breakable,colback=teal!4,colframe=teal!35!black,title={Before you move on}]
What does \ml{കടുക്} mean? (\textquotedblleft{}Mustard seeds\textquotedblright{}.) Say it once more.
\end{tcolorbox}
\section[\texorpdfstring{\ml{മസാല} (masāla)}{മസാല (masāla)}]{\ml{മസാല} (masāla) — spice mix, masala}

\label{lesson:ML-C209-masala}



\begin{quote}
Before the new one: say the Malayalam for curry leaves, then the Malayalam for mustard seeds.
\end{quote}

\subsection*{You'll want to know: \ml{മസാല}}
\textbf{\ml{മസാല}} — \emph{masāla} — \textquotedblleft{}spice mix, masala\textquotedblright{}.

\begin{grammarlens}[title={What you've built}]
One more word for this chapter.
\end{grammarlens}

\subsection*{Your turn}
\begin{itemize}
\raggedright
  \item \emph{Say it:} \emph{masāla}
  \item \emph{Say it:} \emph{masāla}, once more
  \item \emph{Say it:} say \emph{masāla}
  \item \emph{Recall:} say the Malayalam for curry leaves, then the Malayalam for mustard seeds, then say \emph{masāla} again
\end{itemize}

\begin{tcolorbox}[breakable,colback=teal!4,colframe=teal!35!black,title={Before you move on}]
What does \ml{മസാല} mean? (\textquotedblleft{}Spice mix, masala\textquotedblright{}.) Say it once more.
\end{tcolorbox}
\section[\texorpdfstring{\ml{ചെമ്മീൻ} (cemmīn)}{ചെമ്മീൻ (cemmīn)}]{\ml{ചെമ്മീൻ} (cemmīn) — a prawn}

\label{lesson:ML-C209-cemmin}



\begin{quote}
Before the new one: say the Malayalam for mustard seeds, then the Malayalam for spice mix.
\end{quote}

\subsection*{You'll want to know: \ml{ചെമ്മീൻ}}
\textbf{\ml{ചെമ്മീൻ}} — \emph{cemmīn} — \textquotedblleft{}a prawn\textquotedblright{}.

\begin{grammarlens}[title={What you've built}]
One more word for this chapter.
\end{grammarlens}

\subsection*{Your turn}
\begin{itemize}
\raggedright
  \item \emph{Say it:} \emph{cemmīn}
  \item \emph{Say it:} \emph{cemmīn}, once more
  \item \emph{Say it:} say \emph{cemmīn}
  \item \emph{Recall:} say the Malayalam for mustard seeds, then the Malayalam for spice mix, then say \emph{cemmīn} again
\end{itemize}

\begin{tcolorbox}[breakable,colback=teal!4,colframe=teal!35!black,title={Before you move on}]
What does \ml{ചെമ്മീൻ} mean? (\textquotedblleft{}A prawn\textquotedblright{}.) Say it once more.
\end{tcolorbox}
\section[\texorpdfstring{\ml{സന്ദേശം} (sandēśaṁ)}{സന്ദേശം (sandēśaṁ)}]{\ml{സന്ദേശം} (sandēśaṁ) — a message}

\label{lesson:ML-C209-sandesam}



\begin{quote}
Before the new one: say the Malayalam for spice mix, then the Malayalam for a prawn.
\end{quote}

\subsection*{You'll want to know: \ml{സന്ദേശം}}
\textbf{\ml{സന്ദേശം}} — \emph{sandēśaṁ} — \textquotedblleft{}a message\textquotedblright{}.

\begin{grammarlens}[title={What you've built}]
One more word for this chapter.
\end{grammarlens}

\subsection*{Your turn}
\begin{itemize}
\raggedright
  \item \emph{Say it:} \emph{sandēśaṁ}
  \item \emph{Say it:} \emph{sandēśaṁ}, once more
  \item \emph{Say it:} say \emph{sandēśaṁ}
  \item \emph{Recall:} say the Malayalam for spice mix, then the Malayalam for a prawn, then say \emph{sandēśaṁ} again
\end{itemize}

\begin{tcolorbox}[breakable,colback=teal!4,colframe=teal!35!black,title={Before you move on}]
What does \ml{സന്ദേശം} mean? (\textquotedblleft{}A message\textquotedblright{}.) Say it once more.
\end{tcolorbox}
